\documentclass[12pt, a4paper, twoside]{article}
% fonte 12, folha A4, impressão anverso (frente)
\raggedbottom
\usepackage[utf8]{inputenc}
% processar caracteres

\usepackage{amsmath, amsfonts, amssymb}
% ambientes, símbolos e fontes matemáticas.

\usepackage{ac/scripts/meuscomandos}
% meus comandos criados ajuda do bing.com e ChatGPT

\usepackage[left = 3cm, top = 3cm, right = 2cm, bottom = 2cm]{geometry}
% margens- document por geometry

\usepackage[brazilian]{babel}
% \usepackage[brazil]{babel}
% conversões tipográficas e linguísticas brasileiras

 \usepackage{csquotes}

\usepackage{pdfpages}
% inserir outras pages .pdf no LaTeX

\usepackage{titlesec}
\usepackage{titletoc}
% formatação dos títulos das seções

\usepackage{fancyhdr}
% custom style, content e posição do cabeçalho e rodapé

\usepackage{colortbl}
% colorir células de tabelas

\usepackage{setspace}
% espaçamento entre linhas

\usepackage{indentfirst}
% recuar 1° linha cada parágrafo

 \usepackage{natbib}
% gerenciar referências bibliográficas
%%%% NÃO FUNCIONA COM pacote BIBLITEX

% \usepackage{abntex2}
 
\usepackage{bookmark}

%

\usepackage{epigraph}
% adicionar epigrafe

\usepackage{graphicx}
% inserir img no LaTeX

\usepackage{array}
% criar e manipular tabelas

% \usepackage[colorlinks=true, allcolors=black]{hyperref}
% criar links clicáveis (refer cruzadas, citações ou URLs. )

% \usepackage[brazilian,hyperpageref]{backref}
% add list pages onde cada ref bibliog foi citada
%%%% NÃO FUNCIONA COM pacote BIBLITEX


% \usepackage{multirow}
% https://www.ctan.org/pkg/multirow
% criar células +1 linha

\usepackage{float} % https://www.ctan.org/pkg/float
% controlar posição de objetos flutuantes (img e tabelas)

\usepackage{booktabs} % https://www.ctan.org/pkg/booktabs
% criar tabelas mais profissionais

\usepackage{enumitem} % https://www.ctan.org/pkg/enumitem
% criar listas enumeradas (em 1,2,3... ou I, II, III, IV...)

% \usepackage[brazilian,hyperpageref]{backref}
% add list pages onde cada ref bibliog foi citada
%%%% NÃO FUNCIONA COM pacote BIBLITEX

\usepackage{epigraph}
% inserir citações ou epígrafes no texto ou antes de um capítulo

\usepackage{subfigure}
% criar figuras com múltiplas +2img que podem ser posicionadas e rotuladas indepen

\usepackage{anyfontsize}
% selecionar qualquer tamanho de img

\usepackage{caption}
% personalizar as legendas das figuras e tabelas.

\usepackage{adjustbox}
% fornece várias macros para ajustar o conteúdo de caixas.
% macros: é um comando que define uma sequênc de instruções p/ LaTeX executar.
% caixa: é um objeto que contém algum conteúdo, como txt, img ou tabela

\usepackage{bm}
% colocar objetos matemáticos em negrito.


\usepackage{placeins}
% controlar a posição dos objetos flutuantes (como figuras e tabelas)

\usepackage{float}
%  criar e personalizar objetos flutuantes (como figuras e tabelas)

\usepackage{etoolbox}
% ajuda a programar em LaTeX + fácil e poderosa. Ele tem muitas funções % que permitem fazer coisas como modificar comandos existentes, testar % % condições, criar listas e variáveis lógicas. 

\usepackage[style=abnt]{biblatex}
\usepackage[style=abnt, backend=biber]{biblatex}
\addbibresource{main.bib}
% Ensure that main.bib exists and contains valid bibliography entries
% Ensure that main.bib exists and contains valid bibliography entries
\numberwithin{equation}{subsection}
% enumerar equação de acordo com a seção que está

\addbibresource{parte-03/00-bibliografia}

% \bibliographystyle{plain}
% ja estou usando abnt (biblatex acima)
% definir o estilo de format das refer bibliogr LaTeX usando o BiBTeX

\raggedbottom
 % https://latexref.xyz/_005craggedbottom.html
% criar layout de pagina mais agradável


%%%%%%%%%%%%%%%%%%%%%%%%%%%%%%%%%

%\newtheorem{<nome>}{<cabeçalho>}[<contador>]
% <nome> é o nome do ambiente que é definido.
% <cabeçalho> é a palavra que será impressa, em negrito, no início do ambiente.
% <contador> é o nome de um contador existente ao qual o ambiente será subordinado.

\newtheorem{teo}{Teorema}[section]
\newtheorem{lema}[teo]{Lema}
\newtheorem{cor}[teo]{Corolário}
\newtheorem{prop}[teo]{Proposição}
\newtheorem{defi}{Definição}
\newtheorem{exem}{Exemplo}


% (\newtheorem{teo}{Teorema}[section]) criou um ambiente chamado teo para escrever teoremas.
% Cada teorema terá um número seguindo o número da seção em que está.
% Por exemplo: Teorema 1.1, Teorema 1.2, etc.
% 
% Você pode usar esse ambiente assim:
% 
% (\begin{teo} Sejam A e B conjuntos não vazios. Então A x B não é vazio. \end{teo}), ou seja,
% 
% Teorema 1.1 Sejam A e B conjuntos não vazios. Então A x B não é vazio.


%%%%%%%%%% CONFIGURAÇÕES DE PÁGINA E FONTES

\pagestyle{fancy}
\fancyhf{}
\renewcommand{\headrulewidth}{0.4pt}
\setlength{\headheight}{16pt}
% C - Centro, L - Esquerda, R - Direita, O - impar, E - par
\fancyhead[RO, LE]{\thepage}
\renewcommand{\sectionmark}[1]{\markboth{#1}{}}

\titlecontents{section}[0cm]{}{\bf\thecontentslabel\ }{}{\titlerule*[.75pc]{.}\contentspage}
\titlecontents{subsection}[0.75cm]{}{\thecontentslabel\ }{}{\titlerule*[.75pc]{.}\contentspage}

\setcounter{secnumdepth}{3}
\setcounter{tocdepth}{3}

\DeclareCaptionFormat{myformat}{ \centering \fontsize{10}{12}\selectfont#1#2#3}
\captionsetup{format=myformat}

\def\blankpage{%
      \clearpage%
      \thispagestyle{empty}%
      \addtocounter{page}{-1}%
      \null%
      \clearpage}

%%%%%%%%%% CONFIGURAÇÕES DE PÁGINA E FONTE

\pagenumbering{arabic}
\setcounter{page}{2}
\onehalfspacing

\setlength{\parindent}{1.5cm}
\setlength{\parskip}{0.2cm}
\setlength{\intextsep}{0.5cm}

\titlespacing*{\section}{0cm}{0cm}{0.5cm}
\titlespacing*{\subsection}{0cm}{0.5cm}{0.5cm}
\titlespacing*{\subsubsection}{0cm}{0.5cm}{0.5cm}
\titlespacing*{\paragraph}{0cm}{0.5cm}{0.5cm}

\titleformat{\paragraph}
{\normalfont\normalsize\bfseries}{\theparagraph}{1em}{}

\pagenumbering{arabic}
\setcounter{page}{3}
    
\fancyhead[RE, LO]{\nouppercase{\emph\leftmark}}

%%%%%%%%%%%%%%%%%%%%%%%%%%%%%%%%%%%%%%%%%
%%%%%%%%%% INÍCIO DO DOCUMENTO %%%%%%%%%%
%%%%%%%%%%%%%%%%%%%%%%%%%%%%%%%%%%%%%%%%%

\begin{document}

%%%%%%%%%% CAPA %%%%%%%%%%
\begin{titlepage}
\begin{center}

%%%%%%%%% Logo UNB %%%%%%%%%
\begin{figure}[h!]
	\centering
		\includegraphics[scale = 0.8]{ac/logos/unb.png}
	\label{fig:unb}
\end{figure}

%%%%%%%%% Texto ABAIXO LOGO %%%%%%%%%
% meuscomandos: \unb e \dep
{\bf \unb \\
\bf \dep }
\vspace{5cm}

%%%%%%%%% TÍTULO %%%%%%%%%
\setcounter{page}{0}
\null
\textbf{\titulo}
\vspace{2.5cm}

%%%%%%%%% AUTOR (DISCENTE) %%%%%%%%%
\vspace{0.2cm}
\textbf{\autor}
\end{center}
\vspace{1.5cm}

%%%%%%%%% ANO (CONCLUSÃO) %%%%%%%%%
% meuscomandos: \ano
\vspace{5cm}
\begin{center}
{\bf{Brasília} \\ }
\bf{\ano}
\end{center}

\end{titlepage}


%%%%%%%%%% ROSTO %%%%%%%%%%
\thispagestyle{empty}

\begin{center}
\textbf{\autor} \\
\vspace{5cm} % força titulo ficar 5cm para baixo do autor
\textbf{\titulo} \\
\vspace{3cm}  % força orientadores ficar 3cm para baixo do titulo
\small
\orien \prof \\ \coorien \coprof
\end{center}


%%%%%%%%% NATUREZA TRABALHO %%%%%%%%%
\vspace*{3cm}
\begin{flushright}
\begin{minipage}{7.5cm}
 \parbox[t]{7.5cm}{\naturezatrabalho}
\end{minipage}
\end{flushright}

%%%%%%%%% ANO (CONCLUSÃO) %%%%%%%%%
\vspace{5cm}

\begin{center}
{\bf{Brasília} \\ }
\bf{\ano}
\end{center}

%%%%%%%%%% ERRATA %%%%%%%%%%
% \section*{Errata}
% \include{parte-00/01-errata}
% \newpage

\section*{Errata}
\input{parte-00/02-errata.tex}
\markboth{Errata}{Errata}
% \markboth{direita}{esquerda} cabeçalho
\newpage

%%%%%%%%%% FOLHA DE APROVAÇÃO %%%%%%%%%%
% \section*{Folha de Aprovação}
% \include{parte-00/01-aprovacao}
% \newpage

% \section*{Folha de Aprovação}
 %%%%%%%%% Ttulo e Autor%%%%%%%%%
\thispagestyle{empty}

\begin{center}
\textbf{\autor} \\
\vspace{3.75cm} % força titulo ficar 5cm para baixo do autor
\Large\textbf{\titulo} \\
\end{center}
 
 %%%%%%%%% NATUREZA TRABALHO %%%%%%%%%
\vspace{2cm}  % força orientadores ficar 3cm para baixo do titulo
\begin{flushright}
\begin{minipage}{7.5cm}
 \parbox[t]{7.5cm}{\naturezatrabalho}
\end{minipage}
\end{flushright}
 
%%%%%%%%%%%%%%%%%%%%%%%%%%%%%%%%%%
 \vspace*{1.5cm}
  \noindent Trabalho aprovado. Brasília, 27 de Agosto de 2002:
  
\vspace{0.7cm}
\assinatura{\textbf{\profvv} \\ \orienvv} 
\assinatura{\textbf{\coprofvv} \\ \coorienvv}
\assinatura{\textbf{\coprofvv} \\  \coorienvv}


  %%%%%%%%% ANO (CONCLUSÃO) %%%%%%%%%
\vspace{1cm}
\begin{center}
{\bf{Brasília} \\ }
\bf{\ano}
\end{center}
\markboth{Folha de Aprovação}{Folha de Aprovação}
% \markboth{direita}{esquerda} cabeçalho
\newpage

%%%%%%%%%% DEDICATÓRIAS %%%%%%%%%%
% \section*{Dedicatórias}
% \include{parte-00/01-dedicatórias}
% \newpage

\section*{Dedicatórias}



% \vspace*{\fill}
 %  \centering
   \noindent
   % \textit{} \vspace*{\fill}
   This work is dedicated to my family and friends who have supported me throughout my academic journey. This work is dedicated to my family and friends who have supported me throughout my academic journey. This work is dedicated to my family and friends who have supported me throughout my academic journey.
  % \vspace*{\fill}
\markboth{Dedicatórias}{Dedicatórias}
% \markboth{direita}{esquerda} cabeçalho
\newpage

%%%%%%%%%% AGRADECIMENTOS %%%%%%%%%%
% \section*{Agradecimentos}
% \include{parte-00/01-agradecimentos}
% \newpage

\section*{Agradecimentos}
% \vspace*{\fill}
 %  \centering
   \noindent
   % \textit{} \vspace*{\fill}
This work is dedicated to my family and friends who have supported me throughout my academic journey. This work is dedicated to my family and friends who have supported me throughout my academic journey. This work is dedicated to my family and friends who have supported me throughout my academic journey. This work is dedicated to my family and friends who have supported me throughout my academic journey. This work is dedicated to my family and friends who have supported me throughout my academic journey. This work is dedicated to my family and friends who have supported me throughout my academic journey.

% \vspace*{\fill}
\markboth{Agradecimentos}{Agradecimentos}
% \markboth{direita}{esquerda} cabeçalho
\newpage

%%%%%%%%%% EPÍGRAFE %%%%%%%%%%
% \section*{Epígrafe}
% \include{parte-00/01-epigrafe}
% \newpage

 \section*{Epígrafe}

\vspace{20cm}

\setlength{\epigraphwidth}{8cm}
\epigraph{"Lorem Ipsum é simplesmente uma simulação de texto da indústria tipográfica e de impressos, e vem sendo utilizado desde o século XVI"}{Lorem Ipsum}

 \markboth{Epígrafe}{Epígrafe}
% \markboth{direita}{esquerda} cabeçalho
\newpage

%%%%%%%%%% RESUMO %%%%%%%%%%
% \section*{Resumo}
% \include{parte-00/02-resumo}
% \newpage

\section*{Resumo}

Lorem Ipsum é simplesmente uma simulação de texto da indústria tipográfica e de impressos, e vem sendo utilizado desde o século XVI, quando um impressor desconhecido pegou uma bandeja de tipos e os embaralhou para fazer um livro de modelos de tipos. Lorem Ipsum sobreviveu não só a cinco séculos, como também ao salto para a editoração eletrônica, permanecendo essencialmente inalterado. Se popularizou na década de 60, quando a Letraset lançou decalques contendo passagens de Lorem Ipsum, e mais recentemente quando passou a ser integrado a softwares de editoração eletrônica como Aldus PageMaker.
É um fato conhecido de todos que um leitor se distrairá com o conteúdo de texto legível de uma página quando estiver examinando sua diagramação. A vantagem de usar Lorem Ipsum é que ele tem uma distribuição normal de letras, ao contrário de "Conteúdo aqui, conteúdo aqui", fazendo com que ele tenha uma aparência similar a de um texto legível. Muitos softwares de publicação e editores de páginas na internet agora usam Lorem Ipsum como texto-modelo padrão, e uma rápida busca por 'lorem ipsum' mostra vários websites ainda em sua fase de construção. Várias versões novas surgiram ao longo dos anos, eventualmente por acidente, e às vezes de propósito (injetando humor, e coisas do gênero).

 \textbf{Palavras-chaves}: Direito \& Matemática. Direito \& Ciência da Computação. Educação Jurídica. Exame da OAB. Perguntas-Respostas. Experimento Computacional.
\markboth{Resumo}{Resumo}
% \markboth{direita}{esquerda} cabeçalho
\newpage

%%%%%%%%%% ABSTRACT %%%%%%%%%%%
% \section{Abstract}
% \include{parte-00/03-abstract}
% \newpage

\section*{Abstract}

Lorem Ipsum é simplesmente uma simulação de texto da indústria tipográfica e de impressos, e vem sendo utilizado desde o século XVI, quando um impressor desconhecido pegou uma bandeja de tipos e os embaralhou para fazer um livro de modelos de tipos. Lorem Ipsum sobreviveu não só a cinco séculos, como também ao salto para a editoração eletrônica, permanecendo essencialmente inalterado. Se popularizou na década de 60, quando a Letraset lançou decalques contendo passagens de Lorem Ipsum, e mais recentemente quando passou a ser integrado a softwares de editoração eletrônica como Aldus PageMaker.
É um fato conhecido de todos que um leitor se distrairá com o conteúdo de texto legível de uma página quando estiver examinando sua diagramação. A vantagem de usar Lorem Ipsum é que ele tem uma distribuição normal de letras, ao contrário de "Conteúdo aqui, conteúdo aqui", fazendo com que ele tenha uma aparência similar a de um texto legível. Muitos softwares de publicação e editores de páginas na internet agora usam Lorem Ipsum como texto-modelo padrão, e uma rápida busca por 'lorem ipsum' mostra vários websites ainda em sua fase de construção. Várias versões novas surgiram ao longo dos anos, eventualmente por acidente, e às vezes de propósito (injetando humor, e coisas do gênero).

 \textbf{Palavras-chaves:} Direito \& Matemática. Direito \& Ciência da Computação. Educação Jurídica. Exame da OAB. Perguntas-Respostas. Experimento Computacional.
\markboth{Abstract}{Abstract}
% \markboth{direita}{esquerda} cabeçalho
\newpage


%%%%%%%%%% LOCAL- GRÁFICOS - IMAGENS %%%%%%%%%%
\graphicspath{ {arquivos/graficos/} }
\graphicspath{{arquivo/imgs/}}

%%% LISTA DE TABELAS 
\listoftables
\newpage

%%% LISTA DE ILUSTRAÇÕES
\listoffigures
\newpage

%%% LISTA DE ABREVIATURAS E SIGLAS

%%% LISTA DE SÍMBOLOS

%%%%%%%%%% SUMÁRIO %%%%%%%%%%
\tableofcontents

%%%%%%%%%% PÁGINA EM BRANCO %%%%%%%%%%
\blankpage
\newpage

%%%%%%%%%% INTRODUÇÃO %%%%%%%%%%
\section{Introdução}
Curabitur pretium tincidunt lacus. Nulla gravida orci a odio. Nullam varius, turpis et commodo pharetra, est eros bibendum elit, nec luctus magna felis sollicitudin mauris. Integer in mauris eu nibh euismod gravida. Duis ac tellus et risus vulputate vehicula. Donec lobortis risus a elit. Etiam tempor. Ut ullamcorper, ligula eu tempor congue, eros est euismod turpis, id tincidunt sapien risus a quam. Maecenas fermentum consequat mi. Donec fermentum. Pellentesque malesuada nulla a mi. Duis sapien sem, aliquet nec, commodo eget, consequat quis, neque. Aliquam faucibus, elit ut dictum aliquet, felis nisl adipiscing sapien, sed malesuada diam lacus eget erat. Cras mollis scelerisque nunc. Nullam arcu. Aliquam consequat. Curabitur augue lorem, dapibus quis, laoreet et, pretium ac, nisi. Aenean magna nisl, mollis quis, molestie eu, feugiat in, orci.

Curabitur pretium tincidunt lacus. Nulla gravida orci a odio. Nullam varius, turpis et commodo pharetra, est eros bibendum elit, nec luctus magna felis sollicitudin mauris. Integer in mauris eu nibh euismod gravida. Duis ac tellus et risus vulputate vehicula. Donec lobortis risus a elit. Etiam tempor. Ut ullamcorper, ligula eu tempor congue, eros est euismod turpis, id tincidunt sapien risus a quam. Maecenas fermentum consequat mi. Donec fermentum. Pellentesque malesuada nulla a mi. Duis sapien sem, aliquet nec, commodo eget, consequat quis, neque. Aliquam faucibus, elit ut dictum aliquet, felis nisl adipiscing sapien, sed malesuada diam lacus eget erat. Cras mollis scelerisque nunc. Nullam arcu. Aliquam consequat. Curabitur augue lorem, dapibus quis, laoreet et, pretium ac, nisi. Aenean magna nisl, mollis quis, molestie eu, feugiat in, orci.

Curabitur pretium tincidunt lacus. Nulla gravida orci a odio. Nullam varius, turpis et commodo pharetra, est eros bibendum elit, nec luctus magna felis sollicitudin mauris. Integer in mauris eu nibh euismod gravida. Duis ac tellus et risus vulputate vehicula. Donec lobortis risus a elit. Etiam tempor. Ut ullamcorper, ligula eu tempor congue, eros est euismod turpis, id tincidunt sapien risus a quam. Maecenas fermentum consequat mi. Donec fermentum. Pellentesque malesuada nulla a mi. Duis sapien sem, aliquet nec, commodo eget, consequat quis, neque. Aliquam faucibus, elit ut dictum aliquet, felis nisl adipiscing sapien, sed malesuada diam lacus eget erat. Cras mollis scelerisque nunc. Nullam arcu. Aliquam consequat. Curabitur augue lorem, dapibus quis, laoreet et, pretium ac, nisi. Aenean magna nisl, mollis quis, molestie eu, feugiat in, orci.


\newpage

%%%%%%%%%% PÁGINA EM BRANCO %%%%%%%%%%
\blankpage
\newpage

%%%%%%%%%% Referencial Teórico %%%%%%%%%%
\section{Revisão literária}
\input{parte-02/03-revisão-literaria}
\newpage

%%%%%%%%%% PÁGINA EM BRANCO %%%%%%%%%%
\blankpage
\newpage

%%%%%%%%%% OBJETIVOS %%%%%%%%%%
\section{Objetivos}

Curabitur pretium tincidunt lacus. Nulla gravida orci a odio. Nullam varius, turpis et commodo pharetra, est eros bibendum elit, nec luctus magna felis sollicitudin mauris. Integer in mauris eu nibh euismod gravida. Duis ac tellus et risus vulputate vehicula. Donec lobortis risus a elit. Etiam tempor. Ut ullamcorper, ligula eu tempor congue, eros est euismod turpis, id tincidunt sapien risus a quam. Maecenas fermentum consequat mi. Donec fermentum. Pellentesque malesuada nulla a mi. Duis sapien sem, aliquet nec, commodo eget, consequat quis, neque. Aliquam faucibus, elit ut dictum aliquet, felis nisl adipiscing sapien, sed malesuada diam lacus eget erat. Cras mollis scelerisque nunc. Nullam arcu. Aliquam consequat. Curabitur augue lorem, dapibus quis, laoreet et, pretium ac, nisi. Aenean magna nisl, mollis quis, molestie eu, feugiat in, orci.


\subsection{objetivo geral}
Curabitur pretium tincidunt lacus. Nulla gravida orci a odio. Nullam varius, turpis et commodo pharetra, est eros bibendum elit, nec luctus magna felis sollicitudin mauris. Integer in mauris eu nibh euismod gravida. Duis ac tellus et risus vulputate vehicula. Donec lobortis risus a elit. Etiam tempor. Ut ullamcorper, ligula eu tempor congue, eros est euismod turpis, id tincidunt sapien risus a quam. Maecenas fermentum consequat mi. Donec fermentum. Pellentesque malesuada nulla a mi. Duis sapien sem, aliquet nec, commodo eget, consequat quis, neque. Aliquam faucibus, elit ut dictum aliquet, felis nisl adipiscing sapien, sed malesuada diam lacus eget erat. Cras mollis scelerisque nunc. Nullam arcu. Aliquam consequat. Curabitur augue lorem, dapibus quis, laoreet et, pretium ac, nisi. Aenean magna nisl, mollis quis, molestie eu, feugiat in, orci.

Curabitur pretium tincidunt lacus. Nulla gravida orci a odio. Nullam varius, turpis et commodo pharetra, est eros bibendum elit, nec luctus magna felis sollicitudin mauris. Integer in mauris eu nibh euismod gravida. Duis ac tellus et risus vulputate vehicula. Donec lobortis risus a elit. Etiam tempor. Ut ullamcorper, ligula eu tempor congue, eros est euismod turpis, id tincidunt sapien risus a quam. Maecenas fermentum consequat mi. Donec fermentum. Pellentesque malesuada nulla a mi. Duis sapien sem, aliquet nec, commodo eget, consequat quis, neque. Aliquam faucibus, elit ut dictum aliquet, felis nisl adipiscing sapien, sed malesuada diam lacus eget erat. Cras mollis scelerisque nunc. Nullam arcu. Aliquam consequat. Curabitur augue lorem, dapibus quis, laoreet et, pretium ac, nisi. Aenean magna nisl, mollis quis, molestie eu, feugiat in, orci.

\subsection{objetivos especificos}
Curabitur pretium tincidunt lacus. Nulla gravida orci a odio. Nullam varius, turpis et commodo pharetra, est eros bibendum elit, nec luctus magna felis sollicitudin mauris. Integer in mauris eu nibh euismod gravida. Duis ac tellus et risus vulputate vehicula. Donec lobortis risus a elit. Etiam tempor. Ut ullamcorper, ligula eu tempor congue, eros est euismod turpis, id tincidunt sapien risus a quam. Maecenas fermentum consequat mi. Donec fermentum. Pellentesque malesuada nulla a mi. Duis sapien sem, aliquet nec, commodo eget, consequat quis, neque. Aliquam faucibus, elit ut dictum aliquet, felis nisl adipiscing sapien, sed malesuada diam lacus eget erat. Cras mollis scelerisque nunc. Nullam arcu. Aliquam consequat. Curabitur augue lorem, dapibus quis, laoreet et, pretium ac, nisi. Aenean magna nisl, mollis quis, molestie eu, feugiat in, orci.

Curabitur pretium tincidunt lacus. Nulla gravida orci a odio. Nullam varius, turpis et commodo pharetra, est eros bibendum elit, nec luctus magna felis sollicitudin mauris. Integer in mauris eu nibh euismod gravida. Duis ac tellus et risus vulputate vehicula. Donec lobortis risus a elit. Etiam tempor. Ut ullamcorper, ligula eu tempor congue, eros est euismod turpis, id tincidunt sapien risus a quam. Maecenas fermentum consequat mi. Donec fermentum. Pellentesque malesuada nulla a mi. Duis sapien sem, aliquet nec, commodo eget, consequat quis, neque. Aliquam faucibus, elit ut dictum aliquet, felis nisl adipiscing sapien, sed malesuada diam lacus eget erat. Cras mollis scelerisque nunc. Nullam arcu. Aliquam consequat. Curabitur augue lorem, dapibus quis, laoreet et, pretium ac, nisi. Aenean magna nisl, mollis quis, molestie eu, feugiat in, orci.

 Aliquam faucibus, elit ut dictum aliquet, felis nisl adipiscing sapien, sed malesuada diam lacus eget erat. Cras mollis scelerisque nunc. Nullam arcu. Aliquam consequat. Curabitur augue lorem, dapibus quis, laoreet et, pretium ac, nisi. Aenean magna nisl, mollis quis, molestie eu, feugiat in, orci.  Aliquam faucibus, elit ut dictum aliquet, felis nisl adipiscing sapien, sed malesuada diam lacus eget erat. Cras mollis scelerisque nunc. Nullam arcu. Aliquam consequat. Curabitur augue lorem, dapibus quis, laoreet et, pretium ac, nisi. Aenean magna nisl, mollis quis, molestie eu, feugiat in, orci.
\newpage

%%%%%%%%%% PÁGINA EM BRANCO %%%%%%%%%%
\blankpage
\newpage

%%%%%%%%%% METODOLOGIA %%%%%%%%%%
\section{Metodologia}
Lorem ipsum dolor sit amet, consectetur adipisicing elit, sed do eiusmod tempor incididunt ut labore et dolore magna aliqua. Ut enim ad minim veniam, quis nostrud exercitation ullamco laboris nisi ut aliquip ex ea commodo consequat. Duis aute irure dolor in reprehenderit in voluptate velit esse cillum dolore eu fugiat nulla pariatur. Excepteur sint occaecat cupidatat non proident, sunt in culpa qui officia deserunt mollit anim id est laborum.

\subsection{Teste}
Curabitur pretium tincidunt lacus. Nulla gravida orci a odio. Nullam varius, turpis et commodo pharetra, est eros bibendum elit, nec luctus magna felis sollicitudin mauris. Integer in mauris eu nibh euismod gravida. Duis ac tellus et risus vulputate vehicula. Donec lobortis risus a elit. Etiam tempor. Ut ullamcorper, ligula eu tempor congue, eros est euismod turpis, id tincidunt sapien risus a quam. Maecenas fermentum consequat mi. Donec fermentum. Pellentesque malesuada nulla a mi. Duis sapien sem, aliquet nec, commodo eget, consequat quis, neque. Aliquam faucibus, elit ut dictum aliquet, felis nisl adipiscing sapien, sed malesuada diam lacus eget erat. Cras mollis scelerisque nunc. Nullam arcu. Aliquam consequat. Curabitur augue lorem, dapibus quis, laoreet et, pretium ac, nisi. Aenean magna nisl, mollis quis, molestie eu, feugiat in, orci. In hac habitasse platea dictumst.
\subsection{Teste}
Fusce convallis, mauris imperdiet gravida bibendum, nisl turpis suscipit mauris, sed placerat ipsum urna sed risus. In convallis tellus a mauris. Curabitur non elit ut libero tristique sodales. Mauris a lacus. Donec mattis semper leo. In hac habitasse platea dictumst. Vivamus facilisis diam at odio. Mauris dictum, nisi eget consequat elementum, lacus ligula molestie metus, non feugiat orci magna ac sem. Donec turpis. Donec vitae metus. Morbi tristique neque eu mauris. Quisque gravida ipsum non sapien. Proin turpis lacus, scelerisque vitae, elementum at, lobortis ac, quam. Aliquam dictum eleifend risus. In hac habitasse platea dictumst. Etiam sit amet diam. Suspendisse odio. Suspendisse nunc. In semper bibendum libero.

\subsubsection{Teste}
Proin nonummy, lacus eget pulvinar lacinia, pede felis dignissim leo, vitae tristique magna lacus sit amet eros. Nullam ornare. Praesent odio ligula, dapibus sed, tincidunt eget, dictum ac, nibh. Nam quis lacus. Nunc eleifend molestie velit. Morbi lobortis quam eu velit. Donec euismod vestibulum massa. Donec non lectus. Aliquam commodo lacus sit amet nulla. Cras dignissim elit et augue. Nullam non diam. Pellentesque habitant morbi tristique senectus et netus et malesuada fames ac turpis egestas. In hac habitasse platea dictumst. Aenean vestibulum. Sed lobortis elit quis lectus. Nunc sed lacus at augue bibendum dapibus.

Aliquam vehicula sem ut pede. Cras purus lectus, egestas eu, vehicula at, imperdiet sed, nibh. Morbi consectetuer luctus felis. Donec vitae nisi. Aliquam tincidunt feugiat elit. Duis sed elit ut turpis ullamcorper feugiat. Praesent pretium, mauris sed fermentum hendrerit, nulla lorem iaculis magna, pulvinar scelerisque urna tellus a justo. Suspendisse pulvinar massa in metus. Duis quis quam. Proin justo. Curabitur ac sapien. Nam erat. Praesent ut quam.

Vivamus commodo, augue et laoreet euismod, sem sapien tempor dolor, ac egestas sem ligula quis lacus. Donec vestibulum tortor ac lacus. Sed posuere vestibulum nisl. Curabitur eleifend fermentum justo. Nullam imperdiet. Integer sit amet mauris imperdiet risus sollicitudin rutrum. Ut vitae turpis. Nulla facilisi. Quisque tortor velit, scelerisque et, facilisis vel, tempor sed, urna. Vivamus nulla elit, vestibulum eget, semper et, scelerisque eget, lacus. Pellentesque viverra purus. Quisque elit. Donec ut dolor.
7
\subsection{Teste}
Duis volutpat elit et erat. In at nulla at nisl condimentum aliquet. Quisque elementum pharetra lacus. Nunc gravida arcu eget nunc. Nulla iaculis egestas magna. Aliquam erat volutpat. Sed pellentesque orci. Etiam lacus lorem, iaculis sit amet, pharetra quis, imperdiet sit amet, lectus. Integer quis elit ac mi aliquam pretium. Nullam mauris orci, porttitor eget, sollicitudin non, vulputate id, risus. Donec varius enim nec sem. Nam aliquam lacinia enim. Quisque eget lorem eu purus dignissim ultricies. Fusce porttitor hendrerit ante. Mauris urna diam, cursus id, mattis eget, tempus sit amet, risus. Curabitur eu felis. Sed eu mi. Nullam lectus mauris, luctus a, mattis ac, tempus non, leo. Cras mi nulla, rhoncus id, laoreet ut, ultricies id, odio.

Donec imperdiet. Vestibulum auctor tortor at orci. Integer semper, nisi eget suscipit eleifend, erat nisl hendrerit justo, eget vestibulum lorem justo ac leo. Integer sem velit, pharetra in, fringilla eu, fermentum id, felis. Vestibulum sed felis. In elit. Praesent et pede vel ante dapibus condimentum. Donec magna. Quisque id risus. Mauris vulputate pellentesque leo. Duis vulputate, ligula at venenatis tincidunt, orci nunc interdum leo, ac egestas elit sem ut lacus. Etiam non diam quis arcu egestas commodo. Curabitur nec massa ac massa gravida condimentum. Aenean id libero. Pellentesque vitae tellus. Fusce lectus est, accumsan ac, bibendum sed, porta eget, augue. Etiam faucibus. Quisque tempus purus eu ante.
\subsubsection{Teste}
Vestibulum sapien nisl, ornare auctor, consectetuer quis, posuere tristique, odio. Fusce ultrices ullamcorper odio. Ut augue nulla, interdum at, adipiscing non, tristique eget, neque. Pellentesque habitant morbi tristique senectus et netus et malesuada fames ac turpis egestas. Ut pede est, condimentum id, scelerisque ac, malesuada non, quam. Proin eu ligula ac sapien suscipit blandit. Suspendisse euismod. Ut accumsan, neque id gravida luctus, arcu pede sodales felis, vel blandit massa arcu eget ligula. Aenean sed turpis. Pellentesque habitant morbi tristique senectus et netus et malesuada fames ac turpis egestas. Donec sem eros, ornare ut, commodo eu, tempor nec, risus. Donec laoreet dapibus ligula. Praesent orci leo, bibendum nec, ornare et, nonummy in, elit. Donec interdum feugiat leo. Vestibulum ante ipsum primis in faucibus orci luctus et ultrices posuere cubilia Curae; Pellentesque feugiat ullamcorper ipsum. Donec convallis tincidunt urna.
\subsubsection{Teste}
Suspendisse et orci et arcu porttitor pellentesque. Sed lacus nunc, fermentum vel, vehicula in, imperdiet eget, urna. Nam consectetuer euismod nunc. Nulla dignissim posuere nulla. Integer iaculis lacinia massa. Nullam sapien augue, condimentum vel, venenatis id, rhoncus pellentesque, sapien. Donec sed ipsum ultrices turpis consectetuer imperdiet. Duis et ipsum ac nisl laoreet commodo. Mauris eu est. Suspendisse id turpis quis orci euismod consequat. Donec tellus mi, luctus sit amet, ultrices a, convallis eu, lorem. Proin faucibus convallis elit. Maecenas rhoncus arcu at arcu. Proin libero. Proin adipiscing. In quis lorem vitae elit consectetuer pretium. Nullam ligula urna, adipiscing nec, iaculis ut, elementum non, turpis. Fusce pulvinar.
\subsubsection{Teste}
Proin nonummy, lacus eget pulvinar lacinia, pede felis dignissim leo, vitae tristique magna lacus sit amet eros. Nullam ornare. Praesent odio ligula, dapibus sed, tincidunt eget, dictum ac, nibh. Nam quis lacus. Nunc eleifend molestie velit. Morbi lobortis quam eu velit. Donec euismod vestibulum massa. Donec non lectus. Aliquam commodo lacus sit amet nulla. Cras dignissim elit et augue. Nullam non diam. Pellentesque habitant morbi tristique senectus et netus et malesuada fames ac turpis egestas. In hac habitasse platea dictumst. Aenean vestibulum. Sed lobortis elit quis lectus. Nunc sed lacus at augue bibendum dapibus.

Aliquam vehicula sem ut pede. Cras purus lectus, egestas eu, vehicula at, imperdiet sed, nibh. Morbi consectetuer luctus felis. Donec vitae nisi. Aliquam tincidunt feugiat elit. Duis sed elit ut turpis ullamcorper feugiat. Praesent pretium, mauris sed fermentum hendrerit, nulla lorem iaculis magna, pulvinar scelerisque urna tellus a justo. Suspendisse pulvinar massa in metus. Duis quis quam. Proin justo. Curabitur ac sapien. Nam erat. Praesent ut quam.

Vivamus commodo, augue et laoreet euismod, sem sapien tempor dolor, ac egestas sem ligula quis lacus. Donec vestibulum tortor ac lacus. Sed posuere vestibulum nisl. Curabitur eleifend fermentum justo. Nullam imperdiet. Integer sit amet mauris imperdiet risus sollicitudin rutrum. Ut vitae turpis. Nulla facilisi. Quisque tortor velit, scelerisque et, facilisis vel, tempor sed, urna. Vivamus nulla elit, vestibulum eget, semper et, scelerisque eget, lacus. Pellentesque viverra purus. Quisque elit. Donec ut dolor.Fusce convallis, mauris imperdiet gravida bibendum, nisl turpis suscipit mauris, sed placerat ipsum urna sed risus. In convallis tellus a mauris. Curabitur non elit ut libero tristique sodales. Mauris a lacus. Donec mattis semper leo. In hac habitasse platea dictumst. Vivamus facilisis diam at odio. Mauris dictum, nisi eget consequat elementum, lacus ligula molestie metus, non feugiat orci magna ac sem. Donec turpis. Donec vitae metus. Morbi tristique neque eu mauris. Quisque gravida ipsum non sapien. Proin turpis lacus, scelerisque vitae, elementum at, lobortis ac, quam. Aliquam dictum eleifend risus. In hac habitasse platea dictumst. Etiam sit amet diam. Suspendisse odio. Suspendisse nunc. In semper bibendum libero.

\subsubsection{Teste}
Proin nonummy, lacus eget pulvinar lacinia, pede felis dignissim leo, vitae tristique magna lacus sit amet eros. Nullam ornare. Praesent odio ligula, dapibus sed, tincidunt eget, dictum ac, nibh. Nam quis lacus. Nunc eleifend molestie velit. Morbi lobortis quam eu velit. Donec euismod vestibulum massa. Donec non lectus. Aliquam commodo lacus sit amet nulla. Cras dignissim elit et augue. Nullam non diam. Pellentesque habitant morbi tristique senectus et netus et malesuada fames ac turpis egestas. In hac habitasse platea dictumst. Aenean vestibulum. Sed lobortis elit quis lectus. Nunc sed lacus at augue bibendum dapibus.

Aliquam vehicula sem ut pede. Cras purus lectus, egestas eu, vehicula at, imperdiet sed, nibh. Morbi consectetuer luctus felis. Donec vitae nisi. Aliquam tincidunt feugiat elit. Duis sed elit ut turpis ullamcorper feugiat. Praesent pretium, mauris sed fermentum hendrerit, nulla lorem iaculis magna, pulvinar scelerisque urna tellus a justo. Suspendisse pulvinar massa in metus. Duis quis quam. Proin justo. Curabitur ac sapien. Nam erat. Praesent ut quam.

Vivamus commodo, augue et laoreet euismod, sem sapien tempor dolor, ac egestas sem ligula quis lacus. Donec vestibulum tortor ac lacus. Sed posuere vestibulum nisl. Curabitur eleifend fermentum justo. Nullam imperdiet. Integer sit amet mauris imperdiet risus sollicitudin rutrum. Ut vitae turpis. Nulla facilisi. Quisque tortor velit, scelerisque et, facilisis vel, tempor sed, urna. Vivamus nulla elit, vestibulum eget, semper et, scelerisque eget, lacus. Pellentesque viverra purus. Quisque elit. Donec ut dolor.Fusce convallis, mauris imperdiet gravida bibendum, nisl turpis suscipit mauris, sed placerat ipsum urna sed risus. In convallis tellus a mauris. Curabitur non elit ut libero tristique sodales. Mauris a lacus. Donec mattis semper leo. In hac habitasse platea dictumst. Vivamus facilisis diam at odio. Mauris dictum, nisi eget consequat elementum, lacus ligula molestie metus, non feugiat orci magna ac sem. Donec turpis. Donec vitae metus. Morbi tristique neque eu mauris. Quisque gravida ipsum non sapien. Proin turpis lacus, scelerisque vitae, elementum at, lobortis ac, quam. Aliquam dictum eleifend risus. In hac habitasse platea dictumst. Etiam sit amet diam. Suspendisse odio. Suspendisse nunc. In semper bibendum libero.

Vivamus commodo, augue et laoreet euismod, sem sapien tempor dolor, ac egestas sem ligula quis lacus. Donec vestibulum tortor ac lacus. Sed posuere vestibulum nisl. Curabitur eleifend fermentum justo. Nullam imperdiet. Integer sit amet mauris imperdiet risus sollicitudin rutrum. Ut vitae turpis. Nulla facilisi. Quisque tortor velit, scelerisque et, facilisis vel, tempor sed, urna. Vivamus nulla elit, vestibulum eget, semper et, scelerisque eget, lacus. Pellentesque viverra purus. Quisque elit. Donec ut dolor.Fusce convallis, mauris imperdiet gravida bibendum, nisl turpis suscipit mauris, sed placerat ipsum urna sed risus. In convallis tellus a mauris. Curabitur non elit ut libero tristique sodales. Mauris a lacus. Donec mattis semper leo. In hac habitasse platea dictumst. Vivamus facilisis diam at odio. Mauris dictum, nisi eget consequat elementum, lacus ligula molestie metus, non feugiat orci magna ac sem. Donec turpis. Donec vitae metus. Morbi tristique neque eu mauris. Quisque gravida ipsum non sapien. Proin turpis lacus, scelerisque vitae, elementum at, lobortis ac, quam. Aliquam dictum eleifend risus. In hac habitasse platea dictumst. Etiam sit amet diam. Suspendisse odio. Suspendisse nunc. In semper bibendum libero.

\newpage  

%%%%%%%%%% PÁGINA EM BRANCO %%%%%%%%%%
\blankpage
\newpage

%%%%%%%%%% ANALISE DOS RESULTADOS %%%%%%%%%%
\section{Analise e Resultados}
Lorem ipsum dolor sit amet, consectetur adipisicing elit, sed do eiusmod tempor incididunt ut labore et dolore magna aliqua. Ut enim ad minim veniam, quis nostrud exercitation ullamco laboris nisi ut aliquip ex ea commodo consequat. Duis aute irure dolor in reprehenderit in voluptate velit esse cillum dolore eu fugiat nulla pariatur. Excepteur sint occaecat cupidatat non proident, sunt in culpa qui officia deserunt mollit anim id est laborum.

\subsection{Teste}
Curabitur pretium tincidunt lacus. Nulla gravida orci a odio. Nullam varius, turpis et commodo pharetra, est eros bibendum elit, nec luctus magna felis sollicitudin mauris. Integer in mauris eu nibh euismod gravida. Duis ac tellus et risus vulputate vehicula. Donec lobortis risus a elit. Etiam tempor. Ut ullamcorper, ligula eu tempor congue, eros est euismod turpis, id tincidunt sapien risus a quam. Maecenas fermentum consequat mi. Donec fermentum. Pellentesque malesuada nulla a mi. Duis sapien sem, aliquet nec, commodo eget, consequat quis, neque. Aliquam faucibus, elit ut dictum aliquet, felis nisl adipiscing sapien, sed malesuada diam lacus eget erat. Cras mollis scelerisque nunc. Nullam arcu. Aliquam consequat. Curabitur augue lorem, dapibus quis, laoreet et, pretium ac, nisi. Aenean magna nisl, mollis quis, molestie eu, feugiat in, orci. In hac habitasse platea dictumst.
\subsection{Teste}
Fusce convallis, mauris imperdiet gravida bibendum, nisl turpis suscipit mauris, sed placerat ipsum urna sed risus. In convallis tellus a mauris. Curabitur non elit ut libero tristique sodales. Mauris a lacus. Donec mattis semper leo. In hac habitasse platea dictumst. Vivamus facilisis diam at odio. Mauris dictum, nisi eget consequat elementum, lacus ligula molestie metus, non feugiat orci magna ac sem. Donec turpis. Donec vitae metus. Morbi tristique neque eu mauris. Quisque gravida ipsum non sapien. Proin turpis lacus, scelerisque vitae, elementum at, lobortis ac, quam. Aliquam dictum eleifend risus. In hac habitasse platea dictumst. Etiam sit amet diam. Suspendisse odio. Suspendisse nunc. In semper bibendum libero.

\subsubsection{Teste}
Proin nonummy, lacus eget pulvinar lacinia, pede felis dignissim leo, vitae tristique magna lacus sit amet eros. Nullam ornare. Praesent odio ligula, dapibus sed, tincidunt eget, dictum ac, nibh. Nam quis lacus. Nunc eleifend molestie velit. Morbi lobortis quam eu velit. Donec euismod vestibulum massa. Donec non lectus. Aliquam commodo lacus sit amet nulla. Cras dignissim elit et augue. Nullam non diam. Pellentesque habitant morbi tristique senectus et netus et malesuada fames ac turpis egestas. In hac habitasse platea dictumst. Aenean vestibulum. Sed lobortis elit quis lectus. Nunc sed lacus at augue bibendum dapibus.

Aliquam vehicula sem ut pede. Cras purus lectus, egestas eu, vehicula at, imperdiet sed, nibh. Morbi consectetuer luctus felis. Donec vitae nisi. Aliquam tincidunt feugiat elit. Duis sed elit ut turpis ullamcorper feugiat. Praesent pretium, mauris sed fermentum hendrerit, nulla lorem iaculis magna, pulvinar scelerisque urna tellus a justo. Suspendisse pulvinar massa in metus. Duis quis quam. Proin justo. Curabitur ac sapien. Nam erat. Praesent ut quam.

Vivamus commodo, augue et laoreet euismod, sem sapien tempor dolor, ac egestas sem ligula quis lacus. Donec vestibulum tortor ac lacus. Sed posuere vestibulum nisl. Curabitur eleifend fermentum justo. Nullam imperdiet. Integer sit amet mauris imperdiet risus sollicitudin rutrum. Ut vitae turpis. Nulla facilisi. Quisque tortor velit, scelerisque et, facilisis vel, tempor sed, urna. Vivamus nulla elit, vestibulum eget, semper et, scelerisque eget, lacus. Pellentesque viverra purus. Quisque elit. Donec ut dolor.
7
\subsection{Teste}
Duis volutpat elit et erat. In at nulla at nisl condimentum aliquet. Quisque elementum pharetra lacus. Nunc gravida arcu eget nunc. Nulla iaculis egestas magna. Aliquam erat volutpat. Sed pellentesque orci. Etiam lacus lorem, iaculis sit amet, pharetra quis, imperdiet sit amet, lectus. Integer quis elit ac mi aliquam pretium. Nullam mauris orci, porttitor eget, sollicitudin non, vulputate id, risus. Donec varius enim nec sem. Nam aliquam lacinia enim. Quisque eget lorem eu purus dignissim ultricies. Fusce porttitor hendrerit ante. Mauris urna diam, cursus id, mattis eget, tempus sit amet, risus. Curabitur eu felis. Sed eu mi. Nullam lectus mauris, luctus a, mattis ac, tempus non, leo. Cras mi nulla, rhoncus id, laoreet ut, ultricies id, odio.

Donec imperdiet. Vestibulum auctor tortor at orci. Integer semper, nisi eget suscipit eleifend, erat nisl hendrerit justo, eget vestibulum lorem justo ac leo. Integer sem velit, pharetra in, fringilla eu, fermentum id, felis. Vestibulum sed felis. In elit. Praesent et pede vel ante dapibus condimentum. Donec magna. Quisque id risus. Mauris vulputate pellentesque leo. Duis vulputate, ligula at venenatis tincidunt, orci nunc interdum leo, ac egestas elit sem ut lacus. Etiam non diam quis arcu egestas commodo. Curabitur nec massa ac massa gravida condimentum. Aenean id libero. Pellentesque vitae tellus. Fusce lectus est, accumsan ac, bibendum sed, porta eget, augue. Etiam faucibus. Quisque tempus purus eu ante.
\subsubsection{Teste}
Vestibulum sapien nisl, ornare auctor, consectetuer quis, posuere tristique, odio. Fusce ultrices ullamcorper odio. Ut augue nulla, interdum at, adipiscing non, tristique eget, neque. Pellentesque habitant morbi tristique senectus et netus et malesuada fames ac turpis egestas. Ut pede est, condimentum id, scelerisque ac, malesuada non, quam. Proin eu ligula ac sapien suscipit blandit. Suspendisse euismod. Ut accumsan, neque id gravida luctus, arcu pede sodales felis, vel blandit massa arcu eget ligula. Aenean sed turpis. Pellentesque habitant morbi tristique senectus et netus et malesuada fames ac turpis egestas. Donec sem eros, ornare ut, commodo eu, tempor nec, risus. Donec laoreet dapibus ligula. Praesent orci leo, bibendum nec, ornare et, nonummy in, elit. Donec interdum feugiat leo. Vestibulum ante ipsum primis in faucibus orci luctus et ultrices posuere cubilia Curae; Pellentesque feugiat ullamcorper ipsum. Donec convallis tincidunt urna.
\subsubsection{Teste}
Suspendisse et orci et arcu porttitor pellentesque. Sed lacus nunc, fermentum vel, vehicula in, imperdiet eget, urna. Nam consectetuer euismod nunc. Nulla dignissim posuere nulla. Integer iaculis lacinia massa. Nullam sapien augue, condimentum vel, venenatis id, rhoncus pellentesque, sapien. Donec sed ipsum ultrices turpis consectetuer imperdiet. Duis et ipsum ac nisl laoreet commodo. Mauris eu est. Suspendisse id turpis quis orci euismod consequat. Donec tellus mi, luctus sit amet, ultrices a, convallis eu, lorem. Proin faucibus convallis elit. Maecenas rhoncus arcu at arcu. Proin libero. Proin adipiscing. In quis lorem vitae elit consectetuer pretium. Nullam ligula urna, adipiscing nec, iaculis ut, elementum non, turpis. Fusce pulvinar.

\subsubsection{Teste}
Proin nonummy, lacus eget pulvinar lacinia, pede felis dignissim leo, vitae tristique magna lacus sit amet eros. Nullam ornare. Praesent odio ligula, dapibus sed, tincidunt eget, dictum ac, nibh. Nam quis lacus. Nunc eleifend molestie velit. Morbi lobortis quam eu velit. Donec euismod vestibulum massa. Donec non lectus. Aliquam commodo lacus sit amet nulla. Cras dignissim elit et augue. Nullam non diam. Pellentesque habitant morbi tristique senectus et netus et malesuada fames ac turpis egestas. In hac habitasse platea dictumst. Aenean vestibulum. Sed lobortis elit quis lectus. Nunc sed lacus at augue bibendum dapibus.

Aliquam vehicula sem ut pede. Cras purus lectus, egestas eu, vehicula at, imperdiet sed, nibh. Morbi consectetuer luctus felis. Donec vitae nisi. Aliquam tincidunt feugiat elit. Duis sed elit ut turpis ullamcorper feugiat. Praesent pretium, mauris sed fermentum hendrerit, nulla lorem iaculis magna, pulvinar scelerisque urna tellus a justo. Suspendisse pulvinar massa in metus. Duis quis quam. Proin justo. Curabitur ac sapien. Nam erat. Praesent ut quam.

Vivamus commodo, augue et laoreet euismod, sem sapien tempor dolor, ac egestas sem ligula quis lacus. Donec vestibulum tortor ac lacus. Sed posuere vestibulum nisl. Curabitur eleifend fermentum justo. Nullam imperdiet. Integer sit amet mauris imperdiet risus sollicitudin rutrum. Ut vitae turpis. Nulla facilisi. Quisque tortor velit, scelerisque et, facilisis vel, tempor sed, urna. Vivamus nulla elit, vestibulum eget, semper et, scelerisque eget, lacus. Pellentesque viverra purus. Quisque elit. Donec ut dolor.Fusce convallis, mauris imperdiet gravida bibendum, nisl turpis suscipit mauris, sed placerat ipsum urna sed risus. In convallis tellus a mauris. Curabitur non elit ut libero tristique sodales. Mauris a lacus. Donec mattis semper leo. In hac habitasse platea dictumst. Vivamus facilisis diam at odio. Mauris dictum, nisi eget consequat elementum, lacus ligula molestie metus, non feugiat orci magna ac sem. Donec turpis. Donec vitae metus. Morbi tristique neque eu mauris. Quisque gravida ipsum non sapien. Proin turpis lacus, scelerisque vitae, elementum at, lobortis ac, quam. Aliquam dictum eleifend risus. In hac habitasse platea dictumst. Etiam sit amet diam. Suspendisse odio. Suspendisse nunc. In semper bibendum libero.

\subsubsection{Teste}
Proin nonummy, lacus eget pulvinar lacinia, pede felis dignissim leo, vitae tristique magna lacus sit amet eros. Nullam ornare. Praesent odio ligula, dapibus sed, tincidunt eget, dictum ac, nibh. Nam quis lacus. Nunc eleifend molestie velit. Morbi lobortis quam eu velit. Donec euismod vestibulum massa. Donec non lectus. Aliquam commodo lacus sit amet nulla. Cras dignissim elit et augue. Nullam non diam. Pellentesque habitant morbi tristique senectus et netus et malesuada fames ac turpis egestas. In hac habitasse platea dictumst. Aenean vestibulum. Sed lobortis elit quis lectus. Nunc sed lacus at augue bibendum dapibus.

Aliquam vehicula sem ut pede. Cras purus lectus, egestas eu, vehicula at, imperdiet sed, nibh. Morbi consectetuer luctus felis. Donec vitae nisi. Aliquam tincidunt feugiat elit. Duis sed elit ut turpis ullamcorper feugiat. Praesent pretium, mauris sed fermentum hendrerit, nulla lorem iaculis magna, pulvinar scelerisque urna tellus a justo. Suspendisse pulvinar massa in metus. Duis quis quam. Proin justo. Curabitur ac sapien. Nam erat. Praesent ut quam.

Vivamus commodo, augue et laoreet euismod, sem sapien tempor dolor, ac egestas sem ligula quis lacus. Donec vestibulum tortor ac lacus. Sed posuere vestibulum nisl. Curabitur eleifend fermentum justo. Nullam imperdiet. Integer sit amet mauris imperdiet risus sollicitudin rutrum. Ut vitae turpis. Nulla facilisi. Quisque tortor velit, scelerisque et, facilisis vel, tempor sed, urna. Vivamus nulla elit, vestibulum eget, semper et, scelerisque eget, lacus. Pellentesque viverra purus. Quisque elit. Donec ut dolor.Fusce convallis, mauris imperdiet gravida bibendum, nisl turpis suscipit mauris, sed placerat ipsum urna sed risus. In convallis tellus a mauris. Curabitur non elit ut libero tristique sodales. Mauris a lacus. Donec mattis semper leo. In hac habitasse platea dictumst. Vivamus facilisis diam at odio. Mauris dictum, nisi eget consequat elementum, lacus ligula molestie metus, non feugiat orci magna ac sem. Donec turpis. Donec vitae metus. Morbi tristique neque eu mauris. Quisque gravida ipsum non sapien. Proin turpis lacus, scelerisque vitae, elementum at, lobortis ac, quam. Aliquam dictum eleifend risus. In hac habitasse platea dictumst. Etiam sit amet diam. Suspendisse odio. Suspendisse nunc. In semper bibendum libero.

Vivamus commodo, augue et laoreet euismod, sem sapien tempor dolor, ac egestas sem ligula quis lacus. Donec vestibulum tortor ac lacus. Sed posuere vestibulum nisl. Curabitur eleifend fermentum justo. Nullam imperdiet. Integer sit amet mauris imperdiet risus sollicitudin rutrum. Ut vitae turpis. Nulla facilisi. Quisque tortor velit, scelerisque et, facilisis vel, tempor sed, urna. Vivamus nulla elit, vestibulum eget, semper et, scelerisque eget, lacus. Pellentesque viverra purus. Quisque elit. Donec ut dolor.Fusce convallis, mauris imperdiet gravida bibendum, nisl turpis suscipit mauris, sed placerat ipsum urna sed risus. In convallis tellus a mauris. Curabitur non elit ut libero tristique sodales. Mauris a lacus. Donec mattis semper leo. In hac habitasse platea dictumst. Vivamus facilisis diam at odio. Mauris dictum, nisi eget consequat elementum, lacus ligula molestie metus, non feugiat orci magna ac sem. Donec turpis. Donec vitae metus. Morbi tristique neque eu mauris. Quisque gravida ipsum non sapien. Proin turpis lacus, scelerisque vitae, elementum at, lobortis ac, quam. Aliquam dictum eleifend risus. In hac habitasse platea dictumst. Etiam sit amet diam. Suspendisse odio. Suspendisse nunc. In semper bibendum libero.



\begin{table}[htbp]
  \centering
  \caption{Estatísticas de nascimentos e óbitos por ano}
  \begin{tabular}{cccc}
    \toprule
    Ano & Nascidos & Óbitos \\
    \midrule
    2000 & 94,746 & 258 \\
    2001 & 92,979 & 285 \\
    2002 & 93,897 & 298 \\
    2003 & 92,081 & 261 \\
    \addlinespace
    2004 & 91,231 & 278 \\
    2005 & 91,556 & 235 \\
    2006 & 87,483 & 239 \\
    2007 & 85,159 & 232 \\
    \addlinespace
    2008 & 87,462 & 197 \\
    2009 & 87,485 & 194 \\
    2010 & 87,476 & 225 \\
    2011 & 89,742 & 186 \\
    \addlinespace
    2012 & 93,274 & 179 \\
    2013 & 94,822 & 193 \\
    2014 & 99,798 & 198 \\
    2015 & 100,672 & 184 \\
    \addlinespace
    2016 & 95,563 & 189 \\
    2017 & 97,520 & 179 \\
    2018 & 98,872 & 212 \\
    2019 & 96,112 & 201 \\
    \addlinespace
    2020 & 92,768 & 148 \\
    2021 & 90,961 & 195 \\
    \bottomrule
  \end{tabular}
\end{table}

\newpage
%%%%%%%%%% CONCLUSÕES %%%%%%%%%%
\section{Concluões}
Lorem ipsum dolor sit amet, consectetur adipisicing elit, sed do eiusmod tempor incididunt ut labore et dolore magna aliqua. Ut enim ad minim veniam, quis nostrud exercitation ullamco laboris nisi ut aliquip ex ea commodo consequat. Duis aute irure dolor in reprehenderit in voluptate velit esse cillum dolore eu fugiat nulla pariatur. Excepteur sint occaecat cupidatat non proident, sunt in culpa qui officia deserunt mollit anim id est laborum.

\subsection{Teste}
Curabitur pretium tincidunt lacus. Nulla gravida orci a odio. Nullam varius, turpis et commodo pharetra, est eros bibendum elit, nec luctus magna felis sollicitudin mauris. Integer in mauris eu nibh euismod gravida. Duis ac tellus et risus vulputate vehicula. Donec lobortis risus a elit. Etiam tempor. Ut ullamcorper, ligula eu tempor congue, eros est euismod turpis, id tincidunt sapien risus a quam. Maecenas fermentum consequat mi. Donec fermentum. Pellentesque malesuada nulla a mi. Duis sapien sem, aliquet nec, commodo eget, consequat quis, neque. Aliquam faucibus, elit ut dictum aliquet, felis nisl adipiscing sapien, sed malesuada diam lacus eget erat. Cras mollis scelerisque nunc. Nullam arcu. Aliquam consequat. Curabitur augue lorem, dapibus quis, laoreet et, pretium ac, nisi. Aenean magna nisl, mollis quis, molestie eu, feugiat in, orci. In hac habitasse platea dictumst.
\subsection{Teste}
Fusce convallis, mauris imperdiet gravida bibendum, nisl turpis suscipit mauris, sed placerat ipsum urna sed risus. In convallis tellus a mauris. Curabitur non elit ut libero tristique sodales. Mauris a lacus. Donec mattis semper leo. In hac habitasse platea dictumst. Vivamus facilisis diam at odio. Mauris dictum, nisi eget consequat elementum, lacus ligula molestie metus, non feugiat orci magna ac sem. Donec turpis. Donec vitae metus. Morbi tristique neque eu mauris. Quisque gravida ipsum non sapien. Proin turpis lacus, scelerisque vitae, elementum at, lobortis ac, quam. Aliquam dictum eleifend risus. In hac habitasse platea dictumst. Etiam sit amet diam. Suspendisse odio. Suspendisse nunc. In semper bibendum libero.

\subsubsection{Teste}
Proin nonummy, lacus eget pulvinar lacinia, pede felis dignissim leo, vitae tristique magna lacus sit amet eros. Nullam ornare. Praesent odio ligula, dapibus sed, tincidunt eget, dictum ac, nibh. Nam quis lacus. Nunc eleifend molestie velit. Morbi lobortis quam eu velit. Donec euismod vestibulum massa. Donec non lectus. Aliquam commodo lacus sit amet nulla. Cras dignissim elit et augue. Nullam non diam. Pellentesque habitant morbi tristique senectus et netus et malesuada fames ac turpis egestas. In hac habitasse platea dictumst. Aenean vestibulum. Sed lobortis elit quis lectus. Nunc sed lacus at augue bibendum dapibus.

Aliquam vehicula sem ut pede. Cras purus lectus, egestas eu, vehicula at, imperdiet sed, nibh. Morbi consectetuer luctus felis. Donec vitae nisi. Aliquam tincidunt feugiat elit. Duis sed elit ut turpis ullamcorper feugiat. Praesent pretium, mauris sed fermentum hendrerit, nulla lorem iaculis magna, pulvinar scelerisque urna tellus a justo. Suspendisse pulvinar massa in metus. Duis quis quam. Proin justo. Curabitur ac sapien. Nam erat. Praesent ut quam.

Vivamus commodo, augue et laoreet euismod, sem sapien tempor dolor, ac egestas sem ligula quis lacus. Donec vestibulum tortor ac lacus. Sed posuere vestibulum nisl. Curabitur eleifend fermentum justo. Nullam imperdiet. Integer sit amet mauris imperdiet risus sollicitudin rutrum. Ut vitae turpis. Nulla facilisi. Quisque tortor velit, scelerisque et, facilisis vel, tempor sed, urna. Vivamus nulla elit, vestibulum eget, semper et, scelerisque eget, lacus. Pellentesque viverra purus. Quisque elit. Donec ut dolor.
7
\subsection{Teste}
Duis volutpat elit et erat. In at nulla at nisl condimentum aliquet. Quisque elementum pharetra lacus. Nunc gravida arcu eget nunc. Nulla iaculis egestas magna. Aliquam erat volutpat. Sed pellentesque orci. Etiam lacus lorem, iaculis sit amet, pharetra quis, imperdiet sit amet, lectus. Integer quis elit ac mi aliquam pretium. Nullam mauris orci, porttitor eget, sollicitudin non, vulputate id, risus. Donec varius enim nec sem. Nam aliquam lacinia enim. Quisque eget lorem eu purus dignissim ultricies. Fusce porttitor hendrerit ante. Mauris urna diam, cursus id, mattis eget, tempus sit amet, risus. Curabitur eu felis. Sed eu mi. Nullam lectus mauris, luctus a, mattis ac, tempus non, leo. Cras mi nulla, rhoncus id, laoreet ut, ultricies id, odio.

Donec imperdiet. Vestibulum auctor tortor at orci. Integer semper, nisi eget suscipit eleifend, erat nisl hendrerit justo, eget vestibulum lorem justo ac leo. Integer sem velit, pharetra in, fringilla eu, fermentum id, felis. Vestibulum sed felis. In elit. Praesent et pede vel ante dapibus condimentum. Donec magna. Quisque id risus. Mauris vulputate pellentesque leo. Duis vulputate, ligula at venenatis tincidunt, orci nunc interdum leo, ac egestas elit sem ut lacus. Etiam non diam quis arcu egestas commodo. Curabitur nec massa ac massa gravida condimentum. Aenean id libero. Pellentesque vitae tellus. Fusce lectus est, accumsan ac, bibendum sed, porta eget, augue. Etiam faucibus. Quisque tempus purus eu ante.
\subsubsection{Teste}
Vestibulum sapien nisl, ornare auctor, consectetuer quis, posuere tristique, odio. Fusce ultrices ullamcorper odio. Ut augue nulla, interdum at, adipiscing non, tristique eget, neque. Pellentesque habitant morbi tristique senectus et netus et malesuada fames ac turpis egestas. Ut pede est, condimentum id, scelerisque ac, malesuada non, quam. Proin eu ligula ac sapien suscipit blandit. Suspendisse euismod. Ut accumsan, neque id gravida luctus, arcu pede sodales felis, vel blandit massa arcu eget ligula. Aenean sed turpis. Pellentesque habitant morbi tristique senectus et netus et malesuada fames ac turpis egestas. Donec sem eros, ornare ut, commodo eu, tempor nec, risus. Donec laoreet dapibus ligula. Praesent orci leo, bibendum nec, ornare et, nonummy in, elit. Donec interdum feugiat leo. Vestibulum ante ipsum primis in faucibus orci luctus et ultrices posuere cubilia Curae; Pellentesque feugiat ullamcorper ipsum. Donec convallis tincidunt urna.
\subsubsection{Teste}
Suspendisse et orci et arcu porttitor pellentesque. Sed lacus nunc, fermentum vel, vehicula in, imperdiet eget, urna. Nam consectetuer euismod nunc. Nulla dignissim posuere nulla. Integer iaculis lacinia massa. Nullam sapien augue, condimentum vel, venenatis id, rhoncus pellentesque, sapien. Donec sed ipsum ultrices turpis consectetuer imperdiet. Duis et ipsum ac nisl laoreet commodo. Mauris eu est. Suspendisse id turpis quis orci euismod consequat. Donec tellus mi, luctus sit amet, ultrices a, convallis eu, lorem. Proin faucibus convallis elit. Maecenas rhoncus arcu at arcu. Proin libero. Proin adipiscing. In quis lorem vitae elit consectetuer pretium. Nullam ligula urna, adipiscing nec, iaculis ut, elementum non, turpis. Fusce pulvinar.
\newpage

%%%%%%%%%% PÁGINA EM BRANCO %%%%%%%%%%
\blankpage
\newpage

%%%%%%%%%% REFERÊNCIAS %%%%%%%%%%
\section{Referências}
% \printbibliography
\nocite{*} % força a citação aparecer no documento
\newpage

%%%%%%%%%% GLOSSÁRIO %%%%%%%%%%
\section{Glossário}
\input{parte-03/01-glosario}
\newpage

%%%%%%%%%% APÊNDICE %%%%%%%%%%
\section*{Apêndice}
\addcontentsline{toc}{section}{\textbf{Apêndice}}
\input{parte-03/02-apendice}

%%%%%%%%%% ANEXO %%%%%%%%%%
\section{Anexo}
\input{parte-03/03-anexo}
\newpage

%%%%%%%%%% ÍNDICE %%%%%%%%%%
\section{Índice}
\input{parte-03/04-indice.tex}
\newpage

%%%%%%%%%% FIM DO DOCUMENTO %%%%%%%%%%
\end{document}